\section{Introducción: doscientos idiomas no son una cifra de un modelo}

Esta clase no trata de «escalar un poco más el Transformer», sino de un problema de sistemas más difícil: cuando los datos disponibles, las tradiciones de escritura, los estándares lingüísticos, los recursos de evaluación y las necesidades reales están muy desequilibrados, ¿cómo se extiende la traducción automática de unos pocos idiomas de muchos recursos a doscientos idiomas, sin perder calidad, seguridad ni utilidad? NLLB (No Language Left Behind) divide este objetivo en una cadena completa: primero entrevista a los usuarios, luego construye conjuntos de evaluación y datos semilla, después extrae corpus mediante minería, entrena modelos multilingües y contiene el sobreajuste, y por último valida el resultado con métricas automáticas, evaluación humana y revisiones de toxicity.

\lecturefigure{slide-01-title.jpg}{No Language Left Behind: de la cobertura de idiomas a una traducción automática centrada en las personas}{Video de la clase de Stanford 00:00:54}

\begin{knowledgebox}{Lectura de la figura: las tres palabras clave del título}
\term{No Language Left Behind} no garantiza que todos los idiomas ya tengan la misma calidad; es una línea de investigación. \term{Scaling} no se refiere solo a la cantidad de parámetros, sino también al número de idiomas, a las direcciones de traducción, al pipeline de datos y a la capacidad de evaluación. \term{Human-Centered} exige que usuarios reales participen en la definición de necesidades, los estándares de datos, la validación humana y los límites de seguridad. Toda la clase se organiza en torno a estos tres niveles.
\end{knowledgebox}

\teachervoice{Al iniciar, Fan comentó que también podría haber hablado de Llama 2, pero que NLLB está igual de relacionado con el tema de IA generativa del curso y que, además, este problema tiene un significado personal para ella: el inglés es su tercer idioma. Esta motivación nos recuerda que los sistemas multilingües no son un benchmark abstracto, sino que determinan a quién puede servir la tecnología.}

\begin{importantbox}{La tesis de sistemas de esta clase}
La calidad de la traducción automática multilingüe no es función de un solo modelo, sino de un sistema acoplado:
\[
Q = F(H,E,S,M,V),
\]
donde $H$ representa las human needs y los mecanismos de participación, $E$ la evaluation infrastructure, $S$ los source data y el mining pipeline, $M$ el model/training y $V$ la validation y la safety. Si falta cualquiera de estos eslabones, aumentar los parámetros o el volumen de datos puede limitarse a amplificar los sesgos existentes.
\end{importantbox}

\subsection{Ruta de lectura}

Esta clase explica primero por qué «cuántos idiomas se admiten» no basta por sí solo para medir el avance; después, siguiendo el orden de la clase, analiza las entrevistas, FLORES-200, NLLB-Seed, WikiMatrix/LASER3/Stopes, la receta de datos, MoE y el curriculum learning; por último, integra las métricas automáticas, el juicio humano, la toxicity y las preguntas y respuestas de la clase (Q\&A) en un marco de evaluación que admite auditoría.

\subsection{Resumen de la sección}

El verdadero objeto de investigación de NLLB es el «sistema de cobertura de idiomas», no un modelo aislado. Cada conjunto de datos, arquitectura y métrica que aparece más adelante debe entenderse dentro del ciclo completo de necesidades, evidencia y seguridad.
